% Part: first-order-logic
% Chapter: completeness
% Section: lindenbaums-lemma

\documentclass[../../../include/open-logic-section]{subfiles}

\begin{document}

\iftag{FOL}
      {\olfileid{fol}{com}{lin}}
      {\olfileid{pl}{com}{lin}}

\olsection{لمّة ليندنبوم}

\begin{explain}
المقصود بهذه اللمّة أن كل مجموعة متسقة من ال!!{sentence}s لها امتداد
من الجمل يجمع إلى الاتساق كونه !!{complete}. وطريق البرهان أن نزيد
!!{sentence} واحدة في كل مرحلة، على شرط ألا يزول الاتساق. نرتب
الزيادات بحيث كل $!A$ يلحق إما $!A$ وإما $\lnot !A$ عند مرحلة ما.
فإذا أخذنا اتحاد المراحل كلها، وجدنا فيه $!A$ أو نفيها~$\lnot !A$،
فكان كاملًا. ويبقى متسقًا لأن كل مرحلة منه مصونة من التناقض.
\end{explain}

\begin{lem}[لمّة ليندنبوم]
\ollabel{lem:lindenbaum} يمكن توسيع كل مجموعة متسقة~$\OLId{greek-variable}{Gamma}$ في
لغة~$\Lang{L}$ إلى مجموعة !!a{complete} ومتسقة~$\OLId{greek-variable}{Gamma}^*$.
\end{lem}

\begin{proof}
ليكن الفرض اتساق $\OLId{greek-variable}{Gamma}$. نرتب $!A_\OLMathNumeral{0}$ و$!A_\OLMathNumeral{1}$ و\dots{} في قائمة
تستوعب ال!!{sentence}s في~$\Lang L$، ثم نضع $\OLId{greek-variable}{Gamma}_\OLMathNumeral{0} = \OLId{greek-variable}{Gamma}$، و
\[
\OLId{greek-variable}{Gamma}_{\OLId{latin-ordinary}{n}+\OLMathNumeral{1}} =
\begin{cases}
\OLId{greek-variable}{Gamma}_\OLId{latin-ordinary}{n} \cup \{ !A_\OLId{latin-ordinary}{n} \} & \textrm{متى كانت $\Gamma_n \cup \{!A_n\}$
  متسقة؛} \\
\OLId{greek-variable}{Gamma}_\OLId{latin-ordinary}{n} \cup \{ \lnot !A_\OLId{latin-ordinary}{n} \} & \textrm{وإلا.}
\end{cases}
\]
ولتكن $\OLId{greek-variable}{Gamma}^* = \bigcup_{\OLId{latin-ordinary}{n} \geq \OLMathNumeral{0}} \OLId{greek-variable}{Gamma}_\OLId{latin-ordinary}{n}$.

نثبت استقراءً اتساق كل $\OLId{greek-variable}{Gamma}_\OLId{latin-ordinary}{n}$. فأما $\OLId{greek-variable}{Gamma}_\OLMathNumeral{0}$ فمتسقة بحكم
تعريفها. وإذا اخترنا $\OLId{greek-variable}{Gamma}_{\OLId{latin-ordinary}{n}+\OLMathNumeral{1}} = \OLId{greek-variable}{Gamma}_\OLId{latin-ordinary}{n} \cup \{!A_\OLId{latin-ordinary}{n}\}$، فإنما
اخترناها لاتساقها. وإلا كان $\OLId{greek-variable}{Gamma}_{\OLId{latin-ordinary}{n}+\OLMathNumeral{1}} = \OLId{greek-variable}{Gamma}_\OLId{latin-ordinary}{n} \cup \{\lnot !A_\OLId{latin-ordinary}{n}\}$،
فلنثبت اتساق $\OLId{greek-variable}{Gamma}_\OLId{latin-ordinary}{n} \cup \{\lnot !A_\OLId{latin-ordinary}{n}\}$. لو لم تكن متسقة، لكانت
\emph{كل من} $\OLId{greek-variable}{Gamma}_\OLId{latin-ordinary}{n} \cup \{!A_\OLId{latin-ordinary}{n}\}$ و$\OLId{greek-variable}{Gamma}_\OLId{latin-ordinary}{n} \cup \{\lnot !A_\OLId{latin-ordinary}{n}\}$
غير متسقة؛ فيلزم عدم اتساق $\OLId{greek-variable}{Gamma}_\OLId{latin-ordinary}{n}$ بمقتضى
\iftag{FOL}{%
  \tagrefs{prfAX/{fol:axd:prv:prop:provability-exhaustive},
    prfSC/{fol:seq:prv:prop:provability-exhaustive},
    prfND/{fol:ntd:prv:prop:provability-exhaustive},
    prfTab/{fol:tab:prv:prop:provability-exhaustive}}}{%
  \tagrefs{prfAX/{pl:axd:prv:prop:provability-exhaustive},
    prfSC/{pl:seq:prv:prop:provability-exhaustive},
    prfND/{pl:ntd:prv:prop:provability-exhaustive},
    prfTab/{pl:tab:prv:prop:provability-exhaustive}}},
خلافًا لفرض الاستقراء.

ونثبت أيضًا أنه لكل~$\OLId{latin-ordinary}{n}$ ولكل $\OLId{latin-ordinary}{i} < \OLId{latin-ordinary}{n}$، يكون $\OLId{greek-variable}{Gamma}_\OLId{latin-ordinary}{i} \subseteq \OLId{greek-variable}{Gamma}_\OLId{latin-ordinary}{n}$،
باستقراء على~$\OLId{latin-ordinary}{n}$. ففي حالة $\OLId{latin-ordinary}{n}=\OLMathNumeral{0}$ ليس ثمة $\OLId{latin-ordinary}{i} < \OLMathNumeral{0}$، فتصدق الدعوى
خلوًا. وافترض الآن صحتها عند~$\OLId{latin-ordinary}{n}$؛ نريد أن $\OLId{latin-ordinary}{i} < \OLId{latin-ordinary}{n}+\OLMathNumeral{1}$ يقتضي
$\OLId{greek-variable}{Gamma}_\OLId{latin-ordinary}{i} \subseteq \OLId{greek-variable}{Gamma}_{\OLId{latin-ordinary}{n}+\OLMathNumeral{1}}$. وبحسب البناء إما
$\OLId{greek-variable}{Gamma}_{\OLId{latin-ordinary}{n}+\OLMathNumeral{1}} = \OLId{greek-variable}{Gamma}_\OLId{latin-ordinary}{n} \cup \{!A_\OLId{latin-ordinary}{n}\}$ وإما $= \OLId{greek-variable}{Gamma}_\OLId{latin-ordinary}{n} \cup \{\lnot
!A_\OLId{latin-ordinary}{n}\}$؛ فيكون $\OLId{greek-variable}{Gamma}_\OLId{latin-ordinary}{n} \subseteq \OLId{greek-variable}{Gamma}_{\OLId{latin-ordinary}{n}+\OLMathNumeral{1}}$. ومتى كان $\OLId{latin-ordinary}{i} < \OLId{latin-ordinary}{n}+\OLMathNumeral{1}$،
كان $\OLId{greek-variable}{Gamma}_\OLId{latin-ordinary}{i} \subseteq \OLId{greek-variable}{Gamma}_\OLId{latin-ordinary}{n}$: بفرض الاستقراء إن كان $\OLId{latin-ordinary}{i} < \OLId{latin-ordinary}{n}$،
وبالانعكاسية $\OLId{greek-variable}{Gamma}_\OLId{latin-ordinary}{n} \subseteq \OLId{greek-variable}{Gamma}_\OLId{latin-ordinary}{n}$ إن كان $\OLId{latin-ordinary}{i} = \OLId{latin-ordinary}{n}$. ثم تفيد
تعدية~$\subseteq$ أن $\OLId{greek-variable}{Gamma}_\OLId{latin-ordinary}{i} \subseteq \OLId{greek-variable}{Gamma}_{\OLId{latin-ordinary}{n}+\OLMathNumeral{1}}$.

بهذا يتبين اتساق $\OLId{greek-variable}{Gamma}^*$. خذ مجموعة منتهية
$\OLId{greek-variable}{Gamma}' \subseteq \OLId{greek-variable}{Gamma}^*$. فإن كان $\OLId{greek-variable}{Gamma}'=\emptyset$، فالاتساق
ظاهر. وإلا فكل $!B \in \OLId{greek-variable}{Gamma}'$ واقع في~$\OLId{greek-variable}{Gamma}_\OLId{latin-ordinary}{i}$ عند فهرس~$\OLId{latin-ordinary}{i}$.
اختر $\OLId{latin-ordinary}{n}$ أكبر هذه الفهارس. وحيث إن
$\OLId{greek-variable}{Gamma}_\OLId{latin-ordinary}{i} \subseteq \OLId{greek-variable}{Gamma}_\OLId{latin-ordinary}{n}$ متى كان $\OLId{latin-ordinary}{i} \le \OLId{latin-ordinary}{n}$، فكل
$!B \in \OLId{greek-variable}{Gamma}'$ يحقق $!B \in \OLId{greek-variable}{Gamma}_\OLId{latin-ordinary}{n}$، فيكون
$\OLId{greek-variable}{Gamma}' \subseteq \OLId{greek-variable}{Gamma}_\OLId{latin-ordinary}{n}$ مع اتساق~$\OLId{greek-variable}{Gamma}_\OLId{latin-ordinary}{n}$. فثبت اتساق كل
مجموعة منتهية $\OLId{greek-variable}{Gamma}' \subseteq \OLId{greek-variable}{Gamma}^*$. وتفيد \iftag{FOL}{%
  \tagrefs{prfAX/{fol:axd:ptn:prop:proves-compact},
    prfSC/{fol:seq:ptn:prop:proves-compact},
    prfND/{fol:ntd:ptn:prop:proves-compact},
    prfTab/{fol:tab:ptn:prop:proves-compact}}}{%
  \tagrefs{prfAX/{pl:axd:ptn:prop:proves-compact},
    prfSC/{pl:seq:ptn:prop:proves-compact},
    prfND/{pl:ntd:ptn:prop:proves-compact},
    prfTab/{pl:tab:ptn:prop:proves-compact}}}، تكون $\OLId{greek-variable}{Gamma}^*$ متسقة.

وأما الاكتمال، فكل !!{sentence} من $\Frm[\OLId{latin-looped}{L}]$ لها موضع في قائمة بناء
~$\OLId{greek-variable}{Gamma}^*$. فإن كان $!A_\OLId{latin-ordinary}{n} \notin \OLId{greek-variable}{Gamma}^*$، لم يمنع إلحاقها إلا
عدم اتساق $\OLId{greek-variable}{Gamma}_\OLId{latin-ordinary}{n} \cup \{!A_\OLId{latin-ordinary}{n}\}$؛ وحينئذ قد ألحقنا
$\lnot !A_\OLId{latin-ordinary}{n} \in \OLId{greek-variable}{Gamma}^*$. فالمجموعة $\OLId{greek-variable}{Gamma}^*$ !!{complete}.
\end{proof}

\end{document}
